\sectionstart
\section{对话办理与状态变更控制}
\subsection{从自然语言到待确认对象}
对话层优先处理明确安全门与当前业务上下文，再选择本地规则、结构化理解、财务方案或普通只读回答。歧义收款人、缺失金额及不明确对象进入澄清，不能由模型凭空补齐。金额更正会创建新转账草案并撤销旧草案，已结算记录不能被聊天改写成未执行状态。

聊天中的“确认”展示确认入口，不直接执行转账。实际确认使用业务对象、精确令牌和明确布尔确认。对话并非完全只读：创建草案、写会话、撤销待办等可以改变服务端状态，但它们与资金记账不同。

\begin{figure}[htbp]\centering
\caption{普通转账从草案到回执的核验链}
\begin{tikzpicture}[node distance=6mm]
\node[box,text width=11.8cm] (a) {当前请求与澄清\quad 不把历史文字当作资金授权};
\node[box,text width=11.8cm,below=of a] (b) {草案生成\quad 登记收款人、金额、余额、现金路径与前置条件};
\node[box,text width=11.8cm,below=of b] (c) {用户独立确认\quad 对象、金额、令牌与有效窗口};
\node[box,text width=11.8cm,below=of c] (d) {写事务内复核\quad 用户归属、版本、过期、最新条件、原交易};
\node[box,text width=11.8cm,below=of d] (e) {首次有效请求：原子记账\qquad 重试：返回原回执};
\foreach \x/\y in {a/b,b/c,c/d,d/e}\draw[arr](\x)--(\y);
\end{tikzpicture}
\tabnote{不通过的对象保持等待、阻断、失效或撤销状态；不能沿图示直通执行。}
\end{figure}

\subsection{普通转账的草案与确认契约}
普通转账草案核验登记收款人、合法金额、单笔不超过 50,000 元、当前已到账余额、转出后现金底线、未来现金路径和工资前置条件。检查通过时为 \code{pending_confirmation}；工资条件未满足时可以显示 \code{waiting_condition}，但该状态不表示其他检查已经通过。其余不可行条件为 \code{blocked}。

草案绑定用户、账户修订号、随机确认令牌和 600 秒有效期。确认要求 \code{confirmation} 严格为布尔值 true、读取同意、精确令牌以及当前账户版本；与生活目标关联的草案还检查生活版本与无待比较场景等条件。最终普通现金检查与记账位于同一写事务中。

\begin{table}[htbp]\centering\small
\caption{普通转账状态的存储与展示口径}
\begin{tabularx}{\textwidth}{@{}L{4.0cm}Y@{}}\toprule
状态 & 含义及后续处理 \\\midrule
\code{pending_confirmation} & 草案生成时可行；确认时仍须复核 \\
\code{waiting_condition} & 工资等条件未满足；满足后重新准备或复核，不自动付款 \\
\code{blocked} & 当前金额、余额、现金路径或条件不满足 \\
\code{expired} / \code{stale} & 窗口过期，或账户、生活计划发生变化；重建有效草案 \\
\code{revoked} & 用户已撤销尚未结算的草案 \\
\code{settled} / \code{executed} & 数据存储记为 settled；已有交易时展示为 executed，查询原回执 \\\bottomrule
\end{tabularx}
\end{table}

部分状态是读取时投影，不是均以同名字段写回数据库。已结算交易优先投影为已执行；不得仅从展示字符串推断底层状态迁移。已记账交易不能通过撤销草案恢复到未执行。

事务以方案 ID 作为规范交易幂等标识；相同用户和方案已有交易时，在核验原确认信息后返回原回执。更换客户端请求键不会因此产生第二笔扣款。首次记账更新余额、两条方向相反的分录、交易、授权、审计与方案状态，并递增账户修订号。本地事务与唯一约束保证的是沙箱账本内的一致性，不能推广为外部支付系统的“恰好一次”承诺。
\source{backend/service.py: _propose_transfer, _present_proposal, confirm, _post, revoke}

\subsection{多步骤任务与部分完成}
复合请求最多处理八个分句，步骤分别标记只读完成、可确认、阻断、需澄清或拒绝。首次整单确认要求任务令牌、账户修订号、600 秒有效期和明确批准的步骤集合，只能批准当前可执行的步骤；挂失另需卡号后四位，但这不等于生产强身份认证。

执行按顺序推进：未批准的步骤记为跳过，成功步骤保留回执，业务失败步骤记为失败，之后获批准的步骤暂停。正常返回的任务终态 \code{executed} 表示整单执行尝试已经结束，可能包含全部跳过或部分失败，不能译为“全部成功”。重复确认返回原结果，跳过和暂停项不能在同一已结束任务中直接补做，应核对原结果后重新规划。

业务失败可以保留此前成功结果；未捕获异常则由外层事务回滚。转账与申购按前步骤更新后的账本核对资金。普通转账的 50,000 元限额不能未经核查推广为所有任务、所有入口或每日统一限额；“剩下的”金额分配还存在独立的额度计算规则。

当前不同变更入口的确认契约并不统一。单项订阅取消依靠显式确认、订阅状态与既有回执处理重试，没有普通转账草案的全部令牌、TTL 和调用方版本字段。卡片及理财任务的执行分支，也未逐一重复规划阶段的所有条件检查。后续应统一“对象、版本、时限、最新条件、幂等”的动作契约；本文不把局部保证写成全局保证。
\source{backend/agent_ops.py: plan, confirm, _execute; backend/service.py: _cancel_subscription}
